\chapter{Introducción}

\section{Descripción general del proyecto}
El presente trabajo final de carrera documenta el diseño, desarrollo e implementación de un sistema de gestión de turnos académicos. En el ámbito universitario, la coordinación de espacios y servicios especializados, como los laboratorios, bibliotecas, requiere de una organización meticulosa para garantizar que tanto alumnos como docentes y público en general puedan acceder a los recursos de manera equitativa y eficiente.

El sistema de \textbf{Turnos Académicos} surge como una solución tecnológica diseñada específicamente para la \textbf{Universidad Nacional de Salta (UNSa)}, con un enfoque inicial y prioritario para el \textbf{Observatorio Alanis}. El Observatorio, al ser un espacio de divulgación científica y académica, recibe constantes solicitudes de visitas, consultas técnicas y uso de instrumentos, lo que hace indispensable contar con una plataforma que centralice y automatice la reserva de citas.

Esta plataforma web permite que la interacción entre la comunidad universitaria y el Observatorio se realice de manera fluida y previsible. A diferencia de los métodos tradicionales de consulta por correo o presenciales, este sistema permite la autogestión por parte del usuario, quien puede visualizar la disponibilidad real de horarios para las distintas actividades ofrecidas desde cualquier dispositivo con acceso a internet.

Desde el punto de vista administrativo, el sistema provee herramientas robustas para que los gestores del Observatorio y otras dependencias universitarias puedan configurar su capacidad de atención (definida como "mesas habilitadas" o puestos de atención), administrar los tipos de \textbf{trámites} o actividades ofrecidas y gestionar un calendario de feriados o días sin actividad académica. Esta granularidad asegura que el sistema sea adaptable a las particularidades del calendario universitario de la UNSa.

\section{Problemática y Justificación}
En las instituciones académicas, la gestión de la demanda para servicios especializados suele realizarse de forma fragmentada. En el caso del Observatorio Alanis, la falta de un sistema centralizado puede derivar en superposición de horarios, dificultades para llevar un registro estadístico de las visitas y una carga administrativa excesiva para el personal encargado de la atención.

La implementación de este sistema se justifica por la necesidad de:
\begin{itemize}
    \item \textbf{Digitalización de la gestión académica:} Modernizar los procesos de reserva de turnos dentro de la UNSa, alineándose con las políticas de transformación digital de la universidad.
    \item \textbf{Optimización del Observatorio:} Permitir una mejor organización de las visitas guiadas y actividades, asegurando que el aforo sea el adecuado para cada jornada.
    \item \textbf{Mejora en la experiencia del usuario:} Facilitar a los estudiantes y al público externo la obtención de un turno sin necesidad de intermediarios o traslados previos.
    \item \textbf{Trazabilidad y Estadísticas:} Generar datos precisos sobre la afluencia al Observatorio, lo cual es vital para la planificación de recursos y la elaboración de informes de impacto institucional.
\end{itemize}

\section{Objetivos del proyecto}

\subsection{Objetivo general}
Desarrollar e implementar una plataforma web de gestión de turnos que centralice la administración de dependencias y trámites, facilitando el acceso a las personas y optimizando los recursos humanos de la organización.

\subsection{Objetivos específicos}
\begin{itemize}
    \item Diseñar e implementar un catálogo de dependencias y trámites configurables, adaptables a la estructura y necesidades de la institución.
    \item Desarrollar un flujo de reserva de turnos intuitivo y responsivo, permitiendo la selección de horarios disponibles en tiempo real y la captura de datos del usuario.
    \item Incorporar un motor de gestión de disponibilidad horaria que organice la capacidad de atención en función de puestos activos, franjas operativas y días no laborables.
    \item Proveer un panel administrativo integral con herramientas para el control de la operatividad diaria, gestión de catálogo de servicios y monitoreo de turnos.
    \item Implementar mecanismos de notificación y generación de comprobantes digitales para acompañar el ciclo de vida del turno.
    \item Documentar el proceso de desarrollo, la arquitectura del sistema y las principales decisiones técnicas en un informe técnico estructurado.
\end{itemize}

\section{Alcance del sistema}
El proyecto abarca el ciclo completo de la gestión de un turno, desde su concepción en la configuración administrativa hasta su conclusión en la atención presencial.

\textbf{Funcionalidades para el Usuario:}
\begin{itemize}
    \item Navegación por el catálogo de dependencias y trámites disponibles.
    \item Selección de fecha y hora mediante un calendario interactivo.
    \item Ingreso de datos personales y validación mínima de seguridad.
    \item Descarga del comprobante del turno en formato PDF.
\end{itemize}

\textbf{Funcionalidades para la Administración:}
\begin{itemize}
    \item Gestión de ABM (Alta, Baja, Modificación) de dependencias, tipos de trámites y usuarios.
    \item Asociación de trámites a dependencias con configuraciones personalizadas.
    \item Configuración de capacidad de atención por franjas horarias.
    \item Gestión de feriados nacionales, provinciales o institucionales.
    \item Monitoreo en tiempo real de los turnos del día.
    \item Reportes básicos de gestión y estadísticas de atención.
\end{itemize}

\section{Tecnologías utilizadas}
El sistema se desarrolla utilizando el ecosistema de PHP y el framework Laravel, junto con tecnologías web modernas para la reactividad y presentación visual.

\textbf{Backend:}
\begin{itemize}
    \item \textbf{PHP:} Lenguaje de programación principal del lado del servidor.
    \item \textbf{Laravel (v8.x):} Framework web principal, aprovechando su arquitectura MVC, sistema de rutas, migraciones, \textit{seeders}, controladores y Eloquent ORM.
\end{itemize}

\textbf{Frontend:}
\begin{itemize}
    \item \textbf{Blade:} Motor de plantillas de Laravel para el renderizado de la estructura general y \textit{layouts} reutilizables.
    \item \textbf{Vue.js 2:} Framework JavaScript progresivo utilizado para brindar reactividad en componentes dinámicos (selector de horarios y carga en tiempo real de trámites sin recargar la página).
    \item \textbf{AdminLTE / CSS:} Plantilla administrativa basada en Bootstrap para garantizar una interfaz responsiva y adaptable a distintos dispositivos.
\end{itemize}

\textbf{Base de datos:}
\begin{itemize}
    \item \textbf{MySQL / MariaDB:} Motor de base de datos relacional gestionado mediante migraciones y modelos relacionales de Eloquent.
\end{itemize}

\textbf{Herramientas de apoyo:}
\begin{itemize}
    \item \textbf{Composer:} Gestor de dependencias de PHP para la integración de librerías de terceros.
    \item \textbf{NPM / Laravel Mix:} Herramientas para la gestión de paquetes JavaScript y la compilación de recursos web (\textit{assets}).
    \item \textbf{Git:} Sistema de control de versiones para el seguimiento del código fuente.
\end{itemize}

\textbf{Paquetes y Librerías principales:}
\begin{itemize}
    \item \textbf{Spatie Laravel-Permission:} Para la gestión de Roles y Permisos (RBAC) de los usuarios del sistema.
    \item \textbf{Snappy (wkhtmltopdf):} Generación de comprobantes de reserva en formato PDF.
    \item \textbf{Laravel Socialite:} Integración de inicio de sesión con cuentas institucionales de Google.
    \item \textbf{Laravel Mailables:} Envío automático de notificaciones por correo electrónico durante el ciclo de vida del turno.
\end{itemize}

\section{Metodología de desarrollo}
El desarrollo del Sistema de Turnos Académicos se llevó a cabo de forma \textbf{iterativa e incremental}, incorporando funcionalidades por módulos y validando cada una de ellas de manera progresiva.

En una primera etapa se relevaron las necesidades específicas del Observatorio Alanis y la estructura organizacional de la UNSa, definiendo el esquema conceptual de datos. Posterior a la base técnica, se desarrollaron los módulos esenciales (reserva pública de turnos y autenticación administrativa). En etapas sucesivas, se integraron funcionalidades avanzadas como la gestión de disponibilidad horaria por mesas habilitadas, el calendario de feriados, la generación de comprobantes PDF y el envío de notificaciones automáticas.

Durante todo el proceso se utilizaron migraciones y \textit{seeders} para estructurar y poblar la base de datos en entornos de prueba, realizando verificaciones funcionales sobre el flujo completo de obtención y atención de turnos.
